% Part:sets-functions-relations
% Chapter: size-of-sets
% Section: non-enumerability-alt

\documentclass[../../../include/open-logic-section]{subfiles}

\begin{document}

\olfileid[tr]{sfr}{siz}{nen-alt}

\olsection{\printtoken{S}{nonenumerable} Kümeler}

\begin{editorial}
  Bu kesim, \olref[enm-alt]{sec}'teki tanımları, yani $\PosInt$'tan örten bir
  fonksiyon yerine $\Nat$ ile bire bir örten bir eşleme gereğini kullanarak
  $\Bin^\omega$ ve $\Pow{\Nat}$'ın sayılamazlığını ispatlar.
\end{editorial}

\begin{explain}
Doğal sayılar kümesi $\Nat$ sonsuzdur. Ayrıca açıkça !!{enumerable}dir.
Fakat dikkat çekici olgu, \emph{!!{nonenumerable}} kümelerin, yani
!!{enumerable} olmayan kümelerin varlığıdır
(bkz. \olref[sfr][siz][enm-alt]{defn:enumerable}).

Bu şaşırtıcı olabilir. Sonuçta $A$'nın !!{nonenumerable} olduğunu söylemek,
bir !!{bijection} $f\colon\Nat\to A$'nın \emph{bulunmadığını}; yani
$\Nat$'ın sonsuz sayıdaki !!{element}ını $A$'ya eşleyen hiçbir fonksiyonun
$A$'nın tamamını tüketmediğini söylemektir. Dolayısıyla $A$
!!{nonenumerable} ise $A$'nın doğal sayılardan ``daha çok'' !!{element}ı
vardır.

Bir kümenin !!{nonenumerable} olduğunu ispatlamak için uygun bir
!!{bijection}un var olamayacağını göstermelisiniz. Bunu yapmanın en iyi yolu,
$A$'nın !!{element}larını numaralandırmaya yönelik her girişimin en az bir
!!{element}ı dışarıda bırakması gerektiğini göstermektir; bu da hiçbir
$f\colon\Nat\to A$ fonksiyonunun !!{surjective} olmadığını gösterir. Bunu
kurmak için genel bir strateji Cantor'un \emph{köşegen yöntemi}ni
kullanmaktır. $A$'nın !!{element}larından oluşan $x_1$, $x_2$, \dots{}
biçiminde bir liste verildiğinde, kuruluşu gereği bu listede bulunamayacak başka bir $A$
!!{element}ı kurarız.

Fakat bütün bunlar en iyi bir örnekle anlaşılır. İlk örneğimiz, bütün sonsuz
$0$ ve $1$ dizgilerinin kümesi $\Bin^\omega$'dır. (`$\Bin$', ikili anlamına
gelir ve onu yalnızca iki elemanlı $\{0,1\}$ kümesi olarak düşünebiliriz.)
\oliflabeldef{sfr:card-arithmetic:card-opps:sec}{\footnote{Daha kesin olarak
$\Bin^\omega$'nın bütün $0$ ve $1$ $\omega$-dizilerinden oluşan, yani
$\funfromto{\omega}{\{0,1\}}$ ile gösterilen küme olduğunu belirtmeliyiz. Fakat
bunun anlamı ancak \olref[sfr][card-arithmetic][card-opps]{sec}'te açıklığa
kavuşacaktır.}{} Şimdilik bu biraz gevşek ifade kafa karışıklığına neden
olmamalıdır.}
\end{explain}

\begin{thm}
\ollabel{thm:nonenum-bin-omega}
$\Bin^\omega$ !!{nonenumerable}dır.
\end{thm}

\begin{proof}
$\Bin^\omega$'nın !!{element}larından oluşan herhangi bir sonsuz listeyi düşünün.
Her $s_n$'nin sonsuz bir $0$ ve $1$ dizgisi olduğu $s_0$, $s_1$, $s_2$,
\dots{} listemiz olsun. $s_n(m)$, bu listedeki $n$'inci dizginin $m$'inci
basamağı olsun. Böylece listemizi, $s_n(m)$'nin $n$'inci satır ile $m$'inci
sütuna yerleştirildiği bir dizi olarak düşünebiliriz:
\[
\begin{array}{c|c|c|c|c|c}
& 0 & 1 & 2 & 3 & \dots \\\hline
0 & \mathbf{s_{0}(0)} & s_{0}(1) & s_{0}(2) & s_0(3) & \dots \\\hline
1 & s_{1}(0)& \mathbf{s_{1}(1)} & s_1(2) & s_1(3) & \dots \\\hline
2 & s_{2}(0)& s_{2}(1) & \mathbf{s_2(2)} & s_2(3) & \dots \\\hline
3 & s_{3}(0)& s_{3}(1) & s_3(2) & \mathbf{s_3(3)} & \dots \\\hline
\vdots & \vdots & \vdots & \vdots & \vdots & \mathbf{\ddots}
\end{array}
\]
Şimdi bu listede bulunmayan sonsuz bir $0$ ve $1$ dizgisi $d$ kuracağız.
Bunu her girdisini, yani bütün $n\in\Nat$ için $d(n)$'yi belirleyerek
yapacağız. Sezgisel olarak yukarıdaki dizinin köşegeni boyunca aşağı doğru
okur (``köşegen yöntemi'' adı buradan gelir), sonra her $1$'i $0$ ve her
$0$'ı~$1$ yaparız. Daha soyut olarak, köşegenin $n$'inci !!{element}ı
$s_n(n)$'nin $1$ ya da $0$ oluşuna göre $d(n)$'yi $0$ ya da $1$ tanımlarız:
\[
d(n) =
\begin{cases}
1 & \text{$s_{n}(n) = 0$ ise}\\
0 & \text{$s_{n}(n) = 1$ ise.}
\end{cases}
\]
$d$ sonsuz bir $0$ ve $1$ dizgisi olduğundan açıkça $d\in\Bin^\omega$'dır.
Fakat $d$'yi her $n\in\Nat$ için $d(n)\neq s_n(n)$ olacak biçimde kurduk.
Yani $d$, $n$'inci girdisinde $s_n$'den farklıdır. Öyleyse her
$n\in\Nat$ için $d\neq s_n$ ve $d$, $s_0$, $s_1$, $s_2$, \dots{} listesinde
bulunamaz.

$\Bin^\omega$'nın !!{element}larından oluşan gelişigüzel bir sonsuz liste
verildiğinde, listenin $\Bin^\omega$'nın bazı !!{element}larını
dışarıda bırakacağını gösterdik. Dolayısıyla $\Bin^\omega$ kümesinin bir
numaralandırması yoktur; yani $\Bin^\omega$ !!{nonenumerable}dır.
\end{proof}

\begin{explain}
Bu ispat yöntemine ``köşegenleştirme'' denir; çünkü $d$'yi tanımlamak için
dizinin köşegenini kullanır. Ancak köşegenleştirme bir dizinin varlığını
gerektirmez. Gerçek bir dizi ve köşegen söz konusu olmadığında bile benzer bir
fikirle bazı kümelerin !!{nonenumerable} olduğunu gösterebiliriz. Aşağıdaki
sonuç bunun nasıl yapılacağını gösterir.
\end{explain}

\begin{thm}
\ollabel{thm:nonenum-pownat}
$\Pow{\Nat}$ !!{enumerable} değildir.
\end{thm}

\begin{proof}
Aynı biçimde ilerleyerek $\Nat$ alt kümelerinden oluşan her listenin bir
$\Nat$ alt kümesini dışarıda bıraktığını gösteririz. $\Nat$ alt kümelerinin
$N_0$, $N_1$, $N_2$, \ldots{} biçiminde bir listesi olduğunu varsayın. $D$
kümesini $n\in D$ ancak ve ancak $n\notin N_n$ olacak biçimde tanımlayın:
\[
D = \Setabs{n \in \Nat}{n \notin N_n}.
\]
Açıkça $D\subseteq\Nat$'tır. Ancak $D$ listede bulunamaz. Gerçekten kuruluşu
gereği $n\in D$ ancak ve ancak $n\notin N_n$ olduğundan her $n\in\Nat$ için
$D\neq N_n$'dir.
\end{proof}

\begin{explain}
Önceki ispat köşegenden söz etmedi. Yine de onu şöyle canlandırırsanız bir
köşegen içerdiğini düşünebilirsiniz: $N_0$, $N_1$, \dots{} kümelerinin bir
dizide yazıldığını; $N_n$'yi $n$'inci satıra, $m$'yi ancak ve ancak $m\in N_n$
ise $m$'inci sütuna yazarak yerleştirdiğimizi düşünün. Örneğin listedeki
ilk dört kümenin $\{0,1,2,\dots\}$, $\{1,3,5,\dots\}$, $\{0,1,4\}$ ve
$\{2,3,4,\dots\}$ olduğunu varsayın; o zaman dizimiz şöyle başlar:
\[
\begin{array}{r@{}rrrrrrr}
  N_0 = \{ & \mathbf{0}, & 1, & 2, & & & & \dots\}\\
  N_1 = \{ &  & \mathbf{1}, &  & 3, &  & 5, & \dots\}\\
  N_2 = \{ & 0, & 1, &  &  & 4\phantom{,} &  & \}\\
  N_3 = \{ &  &  & 2, & \mathbf{3}, & 4, & & \dots\}\\
  &\vdots & & & & & & \ddots\phantom{\}}
  \end{array}
\]
O zaman $D$, köşegen boyunca aşağı inip $n$ ancak ve ancak köşegende
\emph{yoksa} $n\in D$ alınarak elde edilen kümedir. Dolayısıyla yukarıdaki
durumda $0$ ile $1$'i dışarıda bırakır, $2$'yi içeri alır, $3$'ü dışarıda
bırakır vb.'yiz.
\end{explain}

\begin{prob}
Bütün $f\colon\Nat\to\Nat$ fonksiyonları kümesinin açık bir köşegen
argümanıyla !!{nonenumerable} olduğunu gösterin. Yani $f_1$, $f_2$, \dots{}
bir fonksiyonlar listesi ve her $f_i\colon\Nat\to\Nat$ olduğunda bu listede
bulunmayan bir $g\colon\Nat\to\Nat$ fonksiyonunun var olduğunu gösterin.
\end{prob}

\end{document}
